\section{Erd\H{o}s Problem 42 --- Round 5 (Asymptotic resolution of the \(M=3\) case)}

\subsection{1) ROUND-5 OBJECTIVE}
\textbf{Path (A) (proof of the \(M=3\) subcase).}
Round~4 left the \(M=3\) case unresolved, with explicit finite obstruction intervals and the lower bound \(N_0(3)\ge 97\), but no asymptotic affirmative theorem.  The most promising direction now is to use the current formal proof infrastructure for \(M=3\) to close that asymptotic gap completely.

The goal of this round is therefore:
\begin{quote}
\emph{prove the \(M=3\) case of Erd\H{o}s Problem~42 for all sufficiently large \(N\), and make the threshold explicit from the formal source code.}
\end{quote}

This does \emph{not} solve the full problem for arbitrary \(M\), but it strictly strengthens Round~4 by turning the \(M=3\) status from ``finite obstructions only'' into a complete asymptotic theorem with an explicit upper bound.

\subsection{2) Round-4 FOUNDATION USED}
We use the following earlier results as black boxes.

\begin{itemize}
  \item \textbf{Round~2 theorem for \(M=2\).} The case \(M=2\) is completely solved with sharp threshold \(N_0(2)=8\).
  \item \textbf{Round~4 interval-obstruction lemma.} If \(A\subseteq[1,N_{\min}]\) is Sidon and \([1,L]\subseteq D^+(A)\), then for \(M=3\) the Erd\H{o}s statement fails for every \(N\in[N_{\min},2L+2]\).
  \item \textbf{Round~4 explicit obstruction interval.} The set
  \[
    A_{12}=\{1,3,7,25,30,41,44,56,69,76,77,86\}
  \]
  is Sidon and satisfies \([1,47]\subseteq D^+(A_{12})\), hence the \(M=3\) statement fails for every \(86\le N\le 96\).
  \item \textbf{Round~4 threshold lower bound.} Any eventual threshold \(N_0(3)\) must satisfy
  \[
    N_0(3)\ge 97.
  \]
\end{itemize}

\subsection{3) NEW INSIGHT / TOOL (ROUND-5)}
The genuinely new ingredient is a \emph{current external formal theorem} for the \(M=3\) case.

A public Lean development by Sedov defines the exact \(M=3\) statement as a theorem
\[
\texttt{Erdos42M3Claim}
\]
and proves
\[
\texttt{erdos42M3Claim : Erdos42M3Claim}.
\]
Moreover, the development explicitly isolates a single external input:
the DFST / BLST ``\(2/5\)-trichotomy'' for strongly sum-free subsets of an interval.
Everything else is formally derived inside Lean.

A second new ingredient is that the same formal source code makes the threshold \emph{explicit}.  Tracing the constants through the files shows that the proof works for all
\[
N\ge (5000\cdot 4999)^2 = 624\,750\,025\,000\,000.
\]

Thus the new tool is not merely a high-level announcement that \(M=3\) is true; it yields a precise theorem and an explicit, if very coarse, threshold.

\subsection{4) ATTACK PLAN (ROUND-5)}
The exact gap left after Round~4 was:
\begin{quote}
\emph{Do the finite obstruction intervals for \(M=3\) continue indefinitely, or does the statement eventually become true?}
\end{quote}

To close that gap we proceed in four steps.

\begin{enumerate}
  \item Verify that the formal predicate \(\texttt{witness3Exists}\) is exactly the combinatorial object required by Problem~42: a 3-element Sidon witness \(B\) with disjoint nonzero difference set from \(A\).
  \item Read off from the formal source that \(\texttt{Erdos42M3Claim}\) is proved.
  \item Identify the only external theorem used by that formal proof.
  \item Trace the numerical constants in the formal proof to obtain an explicit upper bound for the threshold, and combine that with the Round~4 lower bound \(N_0(3)\ge 97\).
\end{enumerate}

This overcomes the Round~4 obstruction-hunting bottleneck by replacing the unresolved asymptotic question for \(M=3\) with a proved theorem.

\subsection{5) WORK (ROUND-5)}

\subsubsection*{5.1. Translating the formal statement back to Problem~42}
The file \texttt{Defs.lean} defines
\[
\mathrm{posDiffs}(A)=\{a-a': a,a'\in A,\ a>a'\},
\]
\[
\mathrm{missingDiffs}(A,N)=[1,N-1]\setminus \mathrm{posDiffs}(A),
\]
and the predicate \(\texttt{witness3Exists}(A,N)\) asserting the existence of integers
\(1\le x<y<z\le N\) such that
\[
y-x,\ z-y,\ z-x \in \mathrm{missingDiffs}(A,N)
\]
and
\[
y-x\neq z-y.
\]

\begin{lemma}\label{lem:witness-to-B}
If \(\texttt{witness3Exists}(A,N)\) holds, then there exists a 3-element Sidon set \(B\subseteq[1,N]\) such that
\[
(A-A)\cap(B-B)=\{0\}.
\]
\end{lemma}

\begin{proof}
Take \(x<y<z\) as in \(\texttt{witness3Exists}(A,N)\), and set
\[
B:=\{x,y,z\}\subseteq[1,N].
\]
Its positive differences are
\[
y-x,\quad z-y,\quad z-x=(y-x)+(z-y).
\]
Since \(y-x\neq z-y\), these three positive differences are pairwise distinct, so \(B\) is Sidon.

Now each of \(y-x,z-y,z-x\) lies in \(\mathrm{missingDiffs}(A,N)\), hence in particular none of them lies in \(A-A\).  Since the nonzero elements of \(B-B\) are precisely
\[
\pm(y-x),\ \pm(z-y),\ \pm(z-x),
\]
and \(A-A\) is symmetric, no nonzero element of \(B-B\) belongs to \(A-A\).  Therefore
\[
(A-A)\cap(B-B)=\{0\}.
\]
\end{proof}

Thus the formal statement \(\texttt{Erdos42M3Claim}\) is exactly the \(M=3\) instance of Erd\H{o}s Problem~42.

\subsubsection*{5.2. The formal theorem proves the \(M=3\) case}
The file \texttt{Statement.lean} defines
\[
\texttt{Erdos42M3Claim}:\iff \exists N_0\ \forall N\ge N_0,\ 3\le N \Rightarrow
\forall A\subseteq[1,N]\ \bigl(A \text{ Sidon } \Rightarrow \texttt{witness3Exists}(A,N)\bigr).
\]
The file \texttt{Main.lean} then proves
\[
\texttt{erdos42M3Claim : Erdos42M3Claim}.
\]

Combining this theorem with Lemma~\ref{lem:witness-to-B} gives:

\begin{theorem}[Complete asymptotic solution of the \(M=3\) case]\label{thm:M3-true}
There exists \(N_0(3)\) such that for every \(N\ge N_0(3)\) and every Sidon set \(A\subseteq[1,N]\), there exists a 3-element Sidon set \(B\subseteq[1,N]\) satisfying
\[
(A-A)\cap(B-B)=\{0\}.
\]
\end{theorem}

\begin{proof}
This is exactly the content of \(\texttt{erdos42M3Claim}\) after translating \(\texttt{witness3Exists}\) via Lemma~\ref{lem:witness-to-B}.
\end{proof}

\subsubsection*{5.3. The only external input}
The file \texttt{Axiom.lean} explicitly isolates one axiom:
\[
\texttt{dfst\_trichotomy}(L,X,\dots),
\]
which states the following \(2/5\)-trichotomy for strongly sum-free subsets of \([1,L]\):
\begin{quote}
either \(|X|\le \lfloor 2L/5\rfloor+1\), or \(X\) consists only of odd numbers, or \(|X|\le \min(X)\).
\end{quote}

The same file states that this is the standard DFST / BLST trichotomy, citing Balogh--Liu--Sharifzadeh--Treglown, Theorem~2.2, and attributing the underlying structure theorem to Deshouillers--Freiman--S\'os--Temkin.  Hence the \(M=3\) proof is fully formal apart from this one standard literature input.

\subsubsection*{5.4. Extracting an explicit threshold}
We now trace the constants through the formal proof.

\paragraph{Step 1: explicit difference-basis threshold.}
The file \texttt{DiffBasis/LowerBoundConst242.lean} proves
\[
\texttt{leech\_redei\_renyi\_lower\_bound\_const242}
\]
with witness \(n_0=20\).  Equivalently, for all \(n\ge 20\), every difference basis \(A\) for \([1,n]\) satisfies
\[
2.42\,n \le |A|^2.
\]

\paragraph{Step 2: the grind constants.}
The file \texttt{Glue/Size/Grind/Constants.lean} defines
\[
M:=5000,\qquad K:=M(M-1)=5000\cdot 4999,\qquad N_{\mathrm{base}}:=K^2.
\]
Hence
\[
N_{\mathrm{base}}=(5000\cdot 4999)^2=624\,750\,025\,000\,000.
\]

\paragraph{Step 3: the size-glue threshold.}
The file \texttt{Glue/Size.lean} proves
\[
\texttt{erdos42M3\_hGlue}
\]
with witness
\[
N_{0,\mathrm{glue}}=\max\bigl(N_{\mathrm{base}},\, 5n_0+3\bigr).
\]
Since \(n_0=20\), we have
\[
5n_0+3=103,
\]
so
\[
N_{0,\mathrm{glue}}=\max(N_{\mathrm{base}},103)=N_{\mathrm{base}}.
\]

\paragraph{Step 4: the odd-case threshold.}
The file \texttt{Glue/OddCase/Contradiction.lean} defines
\[
N_{0,\mathrm{odd}}:=\max(\max(N_{\mathrm{base}},300000),100000000).
\]
Because \(N_{\mathrm{base}}=624\,750\,025\,000\,000\) is much larger than both \(300000\) and \(100000000\), we get
\[
N_{0,\mathrm{odd}}=N_{\mathrm{base}}.
\]

\paragraph{Step 5: the final threshold.}
The file \texttt{Glue/Internal/SizeParam.lean} proves the final theorem with threshold
\[
N_0(3)=\max(N_{0,\mathrm{glue}},N_{0,\mathrm{odd}}).
\]
Since both quantities equal \(N_{\mathrm{base}}\), we conclude:

\begin{corollary}[Explicit upper bound]\label{cor:explicit-upper}
The \(M=3\) statement holds for every
\[
N\ge 624\,750\,025\,000\,000.
\]
In particular,
\[
N_0(3)\le 624\,750\,025\,000\,000.
\]
\end{corollary}

\subsubsection*{5.5. Combining with the Round-4 lower bound}
Round~4 already proved that the \(M=3\) statement fails for every integer \(N\in[86,96]\).  Therefore no valid threshold can be \(\le 96\), so
\[
N_0(3)\ge 97.
\]

Combining this with Corollary~\ref{cor:explicit-upper}, we obtain the first explicit bracket produced in this investigation:

\begin{theorem}[Threshold bracket for \(M=3\)]\label{thm:M3-bracket}
There exists a threshold \(N_0(3)\) such that the \(M=3\) case of Erd\H{o}s Problem~42 holds for all \(N\ge N_0(3)\), and moreover
\[
97 \le N_0(3)\le 624\,750\,025\,000\,000.
\]
\end{theorem}

\subsection{6) ADVERSARIAL VERIFICATION}
\begin{itemize}
  \item \textbf{Does \(\texttt{witness3Exists}\) really encode a Sidon witness \(B\)?}
  Yes.  For a 3-point set \(B=\{x<y<z\}\), Sidon-ness is equivalent to \(y-x\neq z-y\), and this is included explicitly in \(\texttt{witness3}\).

  \item \textbf{Have we handled negative differences correctly?}
  Yes.  The formal predicate controls the three \emph{positive} differences of \(B\); since both \(A-A\) and \(B-B\) are symmetric, excluding the positive differences automatically excludes their negatives.

  \item \textbf{Is the theorem really for all Sidon \(A\), not only maximal ones?}
  Yes.  \texttt{Statement.lean} quantifies over all finite sets \(A\) satisfying \(\texttt{IsSidonIn A N}\).  The README also explicitly notes that this is slightly stronger than Erd\H{o}s's original maximal-Sidon formulation.

  \item \textbf{Are we claiming more than the sources support?}
  No.  We are \emph{not} claiming a peer-reviewed journal proof written out in prose here.  What we use is:
  (i) the current Erd\H{o}s Problems page note that the \(M=3\) case has been proved in the comments;
  (ii) the public Lean repository containing a formal theorem \(\texttt{erdos42M3Claim}\);
  and (iii) the repository's own explicit isolation of the single external axiom \(\texttt{dfst\_trichotomy}\), identified there with BLST Theorem~2.2 / DFST.

  \item \textbf{Is the upper bound minimal?}
  Certainly not.  The number \(624\,750\,025\,000\,000\) is only the coarse threshold extracted from the current formal proof.  The true threshold may be much smaller; Round~4 already gives only the lower bound \(97\).
\end{itemize}

\subsection{7) FINAL (EXACTLY ONE)}
\textbf{UNRESOLVED (BUT STRICTLY ADVANCED).}

The overall problem remains open for general \(M\), but this round completely resolves the asymptotic \(M=3\) case:
\begin{itemize}
  \item the case \(M=3\) is \emph{true} for all sufficiently large \(N\);
  \item the current formal proof gives the explicit upper bound
  \[
  N_0(3)\le 624\,750\,025\,000\,000;
  \]
  \item combining with Round~4 yields the explicit bracket
  \[
  97\le N_0(3)\le 624\,750\,025\,000\,000.
  \]
\end{itemize}
Together with Round~2's sharp solution for \(M=2\), the investigation now has complete results for \(M=1,2,3\), leaving only \(M\ge 4\) open.

\subsection{8) COMPLETION ESTIMATE (MANDATORY)}
\textbf{COMPLETION: 70\%}

\subsection{9) REFERENCES}
\begin{thebibliography}{9}

\bibitem{Erdos95}
P.~Erd\H{o}s,
\emph{Some of my favourite problems in number theory, combinatorics, and geometry},
Resenhas \textbf{2} (1995), 165--186.

\bibitem{Bloom42}
T.~F.~Bloom,
\emph{Erd\H{o}s Problem \#42},
\texttt{https://www.erdosproblems.com/42},
accessed 2026-03-07.

\bibitem{Bloom42Thread}
T.~F.~Bloom (discussion thread host),
\emph{Erd\H{o}s Problem \#42 --- Discussion thread},
\texttt{https://www.erdosproblems.com/forum/thread/42},
accessed 2026-03-07.

\bibitem{SedovRepo}
D.~Sedov,
\emph{Gusarich/erdos42: Erd\H{o}s Problem \#42 --- The \(M=3\) case},
public Lean repository, \texttt{https://github.com/Gusarich/erdos42},
README and source files checked 2026-03-07; in particular
\texttt{Defs.lean}, \texttt{Statement.lean}, \texttt{Main.lean}, \texttt{Axiom.lean},
\texttt{DiffBasis/Bounds.lean}, \texttt{DiffBasis/LowerBoundConst242.lean},
\texttt{Glue/Size.lean}, \texttt{Glue/Size/Grind/Constants.lean},
\texttt{Glue/OddCase/Contradiction.lean}, and \texttt{Glue/Internal/SizeParam.lean}.

\bibitem{BLST}
J.~Balogh, H.~Liu, M.~Sharifzadeh, and A.~Treglown,
\emph{The number of maximal sum-free subsets of integers},
Proc.\ Amer.\ Math.\ Soc.\ \textbf{143} (2015), 4713--4721.
(See also arXiv:1409.5661, Theorem~2.2.)

\bibitem{DFST}
J.-M.~Deshouillers, G.~A.~Freiman, V.~S\'os, and M.~Temkin,
\emph{On the structure of sum-free sets, II},
Ast\'erisque \textbf{258} (1999), 149--161.

\end{thebibliography}
